\mxunit{premiere-c05}{Probabilités conditionnelles et indépendance}\label{C05}
\begin{keypoint}[Le parcours du chapitre]
\textbf{Objectifs.} Changer d’univers de référence ; construire un arbre ; employer une partition ; distinguer incompatibilité et indépendance ; traiter jusqu’à quatre Bernoulli.

\textbf{Prérequis.} Fractions, proportions, probabilités sur un univers fini, intersection et réunion.

\textbf{Repère annuel.} 17–26 novembre : 8 h\end{keypoint}
\section*{Diagnostic : six questions, 10 minutes}\label{C05-D}
Répondre sans consulter le cours. Indiquer les questions pour lesquelles une aide est nécessaire. Ce diagnostic ne donne pas lieu à une note.
\begin{enumerate}
\item Sur 40 élèves, 12 pratiquent un sport donné. Proportion ?
\item Si $P(A)=\num{0,3}$, donner $P(\overline A)$.
\item Exprimer « A et B » en notation ensembliste.
\item Exprimer « A ou B » avec le ou inclusif.
\item Calculer $\num{0,4}\times\num{0,25}$.
\item Un dé équilibré est lancé. Donner $P(\ens{2;4;6})$.
\end{enumerate}
\paragraph{Après le diagnostic.} Reprendre les notions signalées avec le professeur ; refaire une question analogue, puis expliquer la règle utilisée.
\paragraph{Oral et logique.} Expliquer à voix haute ce que signifient « sachant », « et » et « au moins ». Sur un arbre, justifier une addition et une multiplication.
\mxnewpage
\section*{Chercher avant de formaliser}\label{C05-RECH}
Construire un tableau de 100 personnes pour lequel A et B sont indépendants, avec 40 A et 30 B. Modifier deux cases en gardant ces marges pour détruire l’indépendance.
\paragraph{Organisation.} Recherche individuelle 5 min, confrontation à deux 5 min, mise en commun 5 min. Garder une trace d’un essai, même non abouti. Une conjecture doit ensuite être justifiée.
\mxnewpage
\section{Conditionner : changer le dénominateur}\label{C05-U1}
\begin{keypoint}[Définition et résultats]
Pour $P(A)>0$, la probabilité de B sachant A est $P_A(B)=\dfrac{P(A\cap B)}{P(A)}$. Dans un tableau d’effectifs équiprobables, elle est le quotient effectif de $A\cap B$ sur effectif de A. En général $P_A(B)\ne P_B(A)$.
\end{keypoint}
\paragraph{Statut et justification.} Définition. En multipliant par $P(A)$, on obtient la règle $P(A\cap B)=P(A)P_A(B)$. La condition $P(A)>0$ est indispensable.
\begin{example}Sur 100 personnes, 40 appartiennent à A, 30 à B et 12 aux deux. $P_A(B)=\dfrac{12}{40}=\num{0,3}$, tandis que $P_B(A)=\dfrac{12}{30}=\num{0,4}$. $P(A\cap B)=\num{0,12}$.\end{example}\label{C05-EX1}
\begin{pitfall}« Parmi les A » fixe le dénominateur. Une probabilité conditionnelle n’est pas une intersection.\end{pitfall}
\begin{application}Sur 60 objets, 20 sont rouges ; parmi les rouges, 5 sont abîmés. Calculer $P(R\cap D)$ et $P_R(D)$.\end{application}\label{C05-AP1}
\section{Arbres, partitions et probabilités totales}\label{C05-U2}
\begin{keypoint}[Définition et résultats]
Une partition $A_1,\ldots,A_k$ est une famille d’événements disjoints dont la réunion est l’univers. Si leurs probabilités sont positives, $P(B)=\sum_iP(A_i)P_{A_i}(B)$. Sur un arbre, les probabilités issues d’un même nœud ont une somme de 1 ; les probabilités d’un chemin se multiplient ; les chemins disjoints se réunissent en additionnant.
\end{keypoint}
\paragraph{Statut et justification.} $B$ est la réunion disjointe des $B\cap A_i$. Par additivité et par définition des probabilités conditionnelles, $P(B)=\sum_iP(B\cap A_i)=\sum_iP(A_i)P_{A_i}(B)$.
\begin{example}Deux machines fournissent \pc{60} et \pc{40} d’une production ; leurs taux de défaut sont \pc{2} et \pc{5}. $P(D)=\num{0,6}\times\num{0,02}+\num{0,4}\times\num{0,05}=\num{0,032}$. Parmi les défauts, la première machine représente $\dfrac{\num{0,012}}{\num{0,032}}=\dfrac{3}{8}$.\end{example}\label{C05-EX2}
\mxinclude{figures/C05-U2}
\begin{pitfall}La moyenne simple des taux n’est correcte que si les poids des branches sont égaux.\end{pitfall}
\begin{application}$P(A)=\dfrac{1}{3}$, $P_A(B)=\dfrac{3}{4}$, $P_{\overline A}(B)=\dfrac{1}{4}$. Calculer $P(B)$.\end{application}\label{C05-AP2}
\section{Indépendance et incompatibilité}\label{C05-U3}
\begin{keypoint}[Définition et résultats]
A et B sont indépendants si $P(A\cap B)=P(A)P(B)$. Si $P(A)>0$, cela équivaut à $P_A(B)=P(B)$. Ils sont incompatibles si $A\cap B=\varnothing$. Deux événements de probabilités strictement positives et incompatibles ne sont pas indépendants.
\end{keypoint}
\paragraph{Statut et justification.} Le critère conditionnel résulte d’une division par $P(A)>0$. Pour des événements incompatibles de probabilités positives, le membre de gauche vaut 0 et le produit de droite est strictement positif.
\begin{example}Sur un dé équilibré, A : résultat pair, B : résultat au moins 5. $P(A)=\dfrac{1}{2}$, $P(B)=\dfrac{1}{3}$, $P(A\cap B)=\dfrac{1}{6}$ : indépendants. A et « résultat impair » sont incompatibles mais non indépendants.\end{example}\label{C05-EX3}
\begin{pitfall}Une indépendance dans un modèle doit être justifiée ou posée explicitement ; « deux événements différents » ne suffit pas.\end{pitfall}
\begin{application}On donne $P(A)=\num{0,4}$, $P(B)=\num{0,5}$, $P(A\cap B)=\num{0,18}$. Sont-ils indépendants ?\end{application}\label{C05-AP3}
\section{Épreuves successives et Bernoulli}\label{C05-U4}
\begin{keypoint}[Définition et résultats]
Une épreuve de Bernoulli a deux résultats : succès, de probabilité p, et échec, de probabilité $1-p$. Dans une répétition de n épreuves indépendantes et identiques, avec $n\leq4$ ici, les probabilités des branches restent p et $1-p$. Un chemin de k succès a pour probabilité $p^k(1-p)^{n-k}$ ; on compte les chemins qui réalisent l’événement.
\end{keypoint}
\paragraph{Statut et justification.} La probabilité d’un chemin est un produit, par indépendance. Les chemins correspondant à un événement sont disjoints : leurs probabilités s’additionnent.
\begin{example}Pour trois essais avec $p=\dfrac{1}{2}$, exactement deux succès correspondent à SSE, SES, ESS. Chaque chemin a probabilité $\dfrac{1}{8}$, donc total $\dfrac{3}{8}$. Au moins un succès a pour probabilité $1-(\dfrac{1}{2})^3=\dfrac{7}{8}$.\end{example}\label{C05-EX4}
\begin{pitfall}Sans remise dans une urne, les probabilités changent généralement après le premier tirage : le modèle indépendant peut être faux.\end{pitfall}
\begin{application}Quatre essais indépendants de probabilité $p=\dfrac{1}{3}$. Calculer la probabilité de quatre échecs puis d’au moins un succès.\end{application}\label{C05-AP4}
\section{Simuler et discuter le modèle}\label{C05-U5}
\begin{keypoint}[Définition et résultats]
Une simulation pseudo-aléatoire produit des fréquences qui fluctuent autour d’une probabilité du modèle. Pour un point uniforme dans le carré $\intff{0}{1}^2$, la proportion vérifiant $y\leq x^2$ estime l’aire sous la parabole ; la proportion vérifiant $x^2+y^2\leq1$, multipliée par 4, estime $\pi$. Les répétitions doivent utiliser de nouveaux tirages.
\end{keypoint}
\paragraph{Statut et justification.} Principe de simulation admis et expérimenté. L’égalité probabilité-aire est une interprétation du tirage uniforme ; l’intégration n’est pas un prérequis. Une simulation ne démontre pas une valeur exacte.
\begin{example}Parmi 10000 points, si 7840 sont dans le quart de disque, l’estimation de $\pi$ vaut $\dfrac{4\times 7840}{10000}=\num{3,136}$. Deux essais de même taille donnent en général des valeurs différentes.\end{example}\label{C05-EX5}
\begin{pitfall}Fixer une graine permet de reproduire un essai, pas de supprimer l’aléa du modèle ni de garantir une précision.\end{pitfall}
\begin{application}Une simulation de 2000 points donne 675 points sous $y=x^2$. Quelle aire estime-t-elle ?\end{application}\label{C05-AP5}
\mxnewpage
\section*{Parcours d’exercices}\label{C05-PARC}
Commencer par les exercices 1 et 2. Choisir ensuite un exercice 3 ou 4 sans aide. Les exercices 5 à 7 sont une réserve de transfert, de recherche ou de programmation : le professeur en désigne un. L’exercice 8 se reprend à distance.
\begingroup
\mxexercisehead{C05-E1}{Un tableau}
Une enquête porte sur 100 usagers de deux services de transport A et B. Parmi eux, 40 utilisent A, 35 utilisent B et 15 utilisent les deux. On choisit au hasard l’un des 100 usagers ; tous ont la même probabilité d’être choisis. Les événements $A$ et $B$ signifient respectivement « l’usager utilise le service A » et « l’usager utilise le service B ».
\begin{enumerate}[label=\arabic*.]
\item Construire et compléter un tableau à double entrée : utilisation ou non de A en lignes, utilisation ou non de B en colonnes. Ajouter les totaux.
\item Calculer $P_A(B)$ et $P_B(A)$. Pour chaque calcul, préciser la population de référence utilisée au dénominateur.
\item Interpréter les deux probabilités par deux phrases distinctes.
\end{enumerate}
\mxexercisehead{C05-E2}{Production}
Une entreprise fabrique des objets dans deux ateliers A et B, qui assurent respectivement $\dfrac14$ et $\dfrac34$ de la production. Parmi les objets produits par A, 4 pour cent sont défectueux ; parmi ceux produits par B, 2 pour cent le sont. On choisit au hasard un objet dans l’ensemble de la production. On note $D$ l’événement « l’objet est défectueux ».
\begin{enumerate}[label=\arabic*.]
\item Représenter la situation par un arbre pondéré et préciser le sens de chaque probabilité du second niveau.
\item Calculer $P(D)$ et interpréter le résultat comme un taux global de défaut.
\item On sait maintenant que l’objet choisi est défectueux. Quelle est la probabilité qu’il ait été fabriqué dans l’atelier A ?
\end{enumerate}
\mxexercisehead{C05-E3}{Indépendance à vérifier}
On lance un dé équilibré à six faces numérotées de 1 à 6. On considère les événements $A$ : « obtenir un nombre pair » et $B$ : « obtenir un nombre supérieur ou égal à 4 ».
\begin{enumerate}[label=\arabic*.]
\item Écrire les ensembles d’issues correspondant à $A$, $B$ et $A\cap B$.
\item Déterminer si les événements $A$ et $B$ sont indépendants. Justifier à l’aide d’un critère numérique.
\item Expliquer pourquoi ils ne sont pas non plus incompatibles.
\end{enumerate}
\mxexercisehead{C05-E4}{Sans remise}
Une urne contient trois boules rouges et deux boules bleues, indiscernables au toucher. On tire au hasard une première boule, puis une seconde sans remettre la première. À chaque tirage, toutes les boules présentes dans l’urne sont équiprobables.
\begin{enumerate}[label=\arabic*.]
\item Construire l’arbre pondéré des deux tirages, en distinguant la couleur obtenue au premier et au second tirage.
\item Calculer la probabilité d’obtenir deux boules rouges.
\item Calculer la probabilité d’obtenir une boule de chaque couleur. Expliquer pourquoi deux chemins de l’arbre doivent être pris en compte.
\end{enumerate}
\mxexercisehead{C05-E5}{Trois essais}
Une expérience comporte un succès de probabilité $\dfrac13$ et un échec. On la répète trois fois de façon indépendante, dans les mêmes conditions.
\begin{enumerate}[label=\arabic*.]
\item Construire un arbre à trois niveaux, ou dresser la liste des huit suites de résultats possibles.
\item Calculer la probabilité d’obtenir exactement un succès. Indiquer les chemins retenus.
\item Calculer la probabilité d’obtenir au moins deux succès en distinguant les cas de deux et de trois succès.
\end{enumerate}
\mxexercisehead{C05-E6}{Détecteur}
Un dispositif contrôle des objets dont 2 pour cent sont défectueux. Il déclenche un signal pour 90 pour cent des objets défectueux et, à tort, pour 5 pour cent des objets sans défaut. On choisit au hasard un objet contrôlé. On note $D$ l’événement « l’objet est défectueux » et $S$ l’événement « le dispositif déclenche un signal ».
\begin{enumerate}[label=\arabic*.]
\item Déterminer la proportion d’objets réellement défectueux parmi les objets signalés. Choisir une représentation et détailler le raisonnement.
\item Donner une valeur exacte, puis un pourcentage arrondi au dixième.
\item Une personne propose directement 90 pour cent comme réponse. Expliquer précisément quelle confusion de probabilités elle commet.
\end{enumerate}
\mxexercisehead{C05-E7}{Ensembles et cardinal}
On travaille dans l’ensemble $E=\ens{1;2;3;4}$. On considère ses parties $A=\ens{1;2}$ et $B=\ens{2;4}$.
\begin{enumerate}[label=\arabic*.]
\item Écrire $A\cap B$, $A\cup B$ et le complémentaire de $A$ dans $E$. Donner le cardinal de chacun.
\item Écrire tous les couples de $A\times B$, puis donner son cardinal.
\item Expliquer pourquoi les couples $\pt{1;2}$ et $\pt{2;1}$ ne peuvent pas être confondus. Le second appartient-il à $A\times B$ ?
\end{enumerate}
\mxexercisehead{C05-E8}{Reprise différée}
Dans une expérience aléatoire, deux événements $A$ et $B$ vérifient\[P(A)=\dfrac12,\qquad P_A(B)=\dfrac14,\qquad P_{\overline A}(B)=\dfrac34.\]
\begin{enumerate}[label=\arabic*.]
\item Calculer $P(B)$ à l’aide d’un arbre ou de la formule des probabilités totales.
\item Calculer $P_B(A)$ en faisant apparaître l’événement commun utilisé au numérateur.
\item Les probabilités $P_A(B)$ et $P_B(A)$ sont-elles égales ici ? Cette égalité est-elle une propriété générale ?
\end{enumerate}
\paragraph{Travail personnel.} Effectuer uniquement la tâche désignée à l’issue de la séance, dans la progression datée : 10–20 min maximum. Les exercices dépendent des notions déjà étudiées ; certaines questions sont donc différées. Reprendre l’exercice 8 une semaine après le chapitre (10 min). Si le temps prévu est dépassé, noter la difficulté et apporter les essais ; ne pas accumuler toute la réserve.
\endgroup
\mxneed{12}\section*{À faire dans le manuel CQFD}
Ces exercices se traitent dans votre manuel. La progression indique la séance ; ils remplacent le travail personnel prévu ce jour-là (10 à 20 minutes au total).
\paragraph{CQFD, exercice 67 p. 336, questions 1, 2 et 3.}\label{C05-M1} Construire un arbre en conservant les événements D et N du manuel. Attention : N désigne un test négatif.
\paragraph{CQFD, exercice 70 p. 336, questions a et b.}\label{C05-M2} Dénombrer les quatre issues équiprobables des deux lancers pour réfuter les égalités proposées.
\begingroup\small
\mxnewpage\section*{QCM 1 : comprendre et choisir}\label{C05-Q1}
Cinq questions pédagogiques à réponse unique. Choisir une réponse par question ; expliquer ensuite un distracteur. Ce format de chapitre ne reproduit pas le format officiel de l’EAM.
\begin{enumerate}
\item $P_A(B)$ a pour dénominateur
\begin{enumerate}[label=\Alph*.,itemsep=2pt]
\item $P(B)$
\item $P(A)$
\item $P(A\cap B)$
\item $1-P(A)$
\end{enumerate}
\item $P(A)=\num{0,4}$, $P_A(B)=\num{0,2}$ : intersection
\begin{enumerate}[label=\Alph*.,itemsep=2pt]
\item $\num{0,08}$
\item $\num{0,6}$
\item $\num{0,5}$
\item $\num{0,2}$
\end{enumerate}
\item Une partition comporte des événements
\begin{enumerate}[label=\Alph*.,itemsep=2pt]
\item égaux
\item tous indépendants
\item de même probabilité
\item disjoints et couvrant l’univers
\end{enumerate}
\item Indépendance équivaut à
\begin{enumerate}[label=\Alph*.,itemsep=2pt]
\item $P(A)=P(B)$
\item $P(A\cap B)=0$
\item $P(A\cap B)=P(A)P(B)$
\item $A=B$
\end{enumerate}
\item Deux événements incompatibles de probabilités positives sont
\begin{enumerate}[label=\Alph*.,itemsep=2pt]
\item indépendants
\item non indépendants
\item égaux
\item complémentaires nécessairement
\end{enumerate}
\end{enumerate}
\endgroup
\begingroup\small
\mxneed{18}\section*{QCM 2 : comprendre et choisir}\label{C05-Q2}
Cinq questions pédagogiques à réponse unique. Choisir une réponse par question ; expliquer ensuite un distracteur. Ce format de chapitre ne reproduit pas le format officiel de l’EAM.
\begin{enumerate}
\item Trois lancers indépendants d’une pièce équilibrée, au moins un pile :
\begin{enumerate}[label=\Alph*.,itemsep=2pt]
\item $\dfrac{1}{8}$
\item $\dfrac{3}{8}$
\item $\dfrac{7}{8}$
\item $\dfrac{1}{2}$
\end{enumerate}
\item Deux essais indépendants avec $p=\dfrac{1}{3}$, deux succès :
\begin{enumerate}[label=\Alph*.,itemsep=2pt]
\item $\dfrac{1}{9}$
\item $\dfrac{2}{3}$
\item $\dfrac{1}{3}$
\item $\dfrac{4}{9}$
\end{enumerate}
\item Les branches issues d’un même nœud ont pour somme
\begin{enumerate}[label=\Alph*.,itemsep=2pt]
\item $0$
\item $2$
\item la probabilité du chemin
\item $1$
\end{enumerate}
\item Une simulation donne une fréquence égale à \num{0,51}. Cela
\begin{enumerate}[label=\Alph*.,itemsep=2pt]
\item prouve $p=\num{0,51}$
\item est une observation fluctuante
\item prouve $p=\dfrac{1}{2}$
\item réfute $p=\dfrac{1}{2}$
\end{enumerate}
\item Au tirage sans remise, après une rouge parmi 3 rouges et 2 bleues, probabilité de rouge :
\begin{enumerate}[label=\Alph*.,itemsep=2pt]
\item $\dfrac{3}{5}$
\item $\dfrac{3}{4}$
\item $\dfrac{1}{2}$
\item $\dfrac{2}{5}$
\end{enumerate}
\end{enumerate}
\subsection*{Vrai ou faux ? Justifier.}\label{C05-VF}
\begin{enumerate}
\item $P_A(B)=P_B(A)$ pour tous événements de probabilités positives.
\item Si A et B sont indépendants, A et $\overline B$ aussi.
\item Deux tirages successifs sont toujours indépendants.
\end{enumerate}
\begin{keypoint}[Bilan personnel]Je repère le bon dénominateur ; je vérifie les sommes aux nœuds ; je justifie l’indépendance ; je compte les chemins sans les confondre.

Pour chaque objectif, indiquer : acquis sans aide ; réussi avec aide ; à retravailler. Écrire un exemple et une erreur à éviter.\end{keypoint}
\endgroup
\ifmxteacher\mxinclude{corriges/C05}\fi
